\section{Reference-Adjusted Evaluation}

We introduce the reference-adjusted evaluation paradigm as a middle ground to the current reference-based and reference-free paradigms. In the \textit{reference-based} setting, the LLM judge is provided with one or more gold human-written reference summaries alongside the candidate summary to be evaluated. In the \textit{reference-free} setting, no reference is provided and the judge evaluates the candidate summary in isolation. Our proposed paradigm, \textit{reference-adjusted evaluation}, occupies a middle ground: the judge is provided with one or more references that are related but non-exact. For example, machine-generated summaries or naturally occurring data that differs in style but matches the content.

The motivation for reference-adjusted evaluation stems from two observations. 
First, LLM judges perform better when provided with a reference, as the 
reference gives the model an axis along which to compare the candidate 
summary. Second, gold references are often unavailable or prohibitively 
expensive to collect, particularly when evaluating on new test domains or 
datasets. Reference-adjusted evaluation addresses this tension by leveraging 
cheaply available related summaries as approximate references.

We experiment with content-related references that share the informational content of the ideal summary but may differ in style or wording. In our experiments, we use machine-generated summaries produced by models other than the one being evaluated. Specifically, we experiment with both 
abstractive and extractive references. The intuition is that what other summarization algorithms choose to include is a useful signal about which content is most salient, even if the precise wording does not match the target style.

To communicate the approximate nature of the reference to the LLM judge, 
we modify the standard reference-based prompt template. Rather than presenting 
the reference as a gold standard, we explicitly condition the judge's 
interpretation with the instruction that the reference matches the content 
of the ideal response but does not necessarily match its style or wording. 
This framing is important: it prevents the judge from penalizing a candidate 
summary for legitimate stylistic differences from the reference, while still 
allowing the reference to anchor the judge's assessment of content coverage 
and factual consistency. The full prompt templates for the reference-based, 
reference-adjusted, and reference-free settings are provided in 
Appendix~\ref{sec:6-appendix-prompts}.